\chapter{Hair Loss in Men}

\section*{Definition and Overview}
Hair loss in men is most commonly due to male pattern baldness (androgenetic alopecia), a genetic and hormone-driven condition in which hair thins and recedes in a characteristic pattern, usually starting at the temples and crown. It is very common and, while harmless to health, can affect confidence and self-esteem. Less common causes include temporary shedding after illness or stress (telogen effluvium), autoimmune patchy hair loss (alopecia areata), scalp infections, and medical conditions such as thyroid disease. Treatment depends on the cause, with options including medicines that slow or reverse loss, and hair transplant surgery.

\section*{Epidemiology}
Male pattern hair loss is the most common cause of hair loss in men. Around 30--50% of men have noticeable balding by age 50, and more than 70% are affected to some degree by age 70. It is more common and earlier in onset in men with a family history. Other causes are less common: alopecia areata affects about 2% of people, and telogen effluvium is a frequent temporary cause after illness, stress, or weight loss. Hair loss has a significant psychological impact for many men, with studies linking it to reduced self-esteem.

\section*{Aetiology and Pathophysiology}
The main cause of hair loss in men is the effect of hormones on genetically susceptible hair follicles.

\begin{itemize}
    \item \textbf{Androgenetic alopecia:} the male hormone testosterone is converted to dihydrotestosterone (DHT) by the enzyme 5-alpha reductase; DHT shrinks genetically susceptible follicles in the scalp, causing hair to become shorter, thinner, and finally stop growing
    \item \textbf{Genetic susceptibility:} inherited genes determine which follicles are sensitive
    \item \textbf{Pattern:} typical recession at the temples and thinning on the crown, progressing over years
    \item \textbf{Telogen effluvium:} stress, illness, surgery, or rapid weight loss shifts many follicles into the resting phase, causing diffuse shedding two to three months later
    \item \textbf{Alopecia areata:} the immune system attacks follicles, causing patches of loss
    \item \textbf{Other causes:} scalp infections, thyroid disease, iron deficiency, and medicines
\end{itemize}

\section*{Clinical Features}
The pattern of hair loss is often diagnostic.

\begin{itemize}
    \item Receding hairline at the temples
    \item Thinning on the crown of the head
    \item Gradual progression over years
    \item Diffuse shedding after illness or stress (telogen effluvium)
    \item Patchy, circular bald spots (alopecia areata)
    \item Itching, scaling, or redness of the scalp with some conditions
    \item Hair loss with weight change, fatigue, or other symptoms suggesting an underlying condition
    \item Bald patches with broken hairs and "exclamation mark" hairs in alopecia areata
\end{itemize}

\section*{Differential Diagnosis}
Other causes of hair loss in men should be considered.

\begin{itemize}
    \item Male pattern hair loss (the most common)
    \item Telogen effluvium (temporary diffuse shedding)
    \item Alopecia areata and its variants
    \item Tinea capitis (fungal scalp infection)
    \item Thyroid disease
    \item Iron deficiency and nutritional deficiency
    \item Medicine-related hair loss
    \item Scarring alopecia from inflammatory scalp conditions
    \item Trichotillomania (hair pulling)
\end{itemize}

\section*{When to Seek Medical Care}
See a doctor if hair loss is sudden, patchy, associated with a rash or scaling, or accompanied by other symptoms, or if you want to discuss treatment options for pattern hair loss.

\textbf{Call 000 (or local emergency services) for:}
\begin{itemize}
    \item This condition is not an emergency; however, sudden hair loss with scalp pain, swelling, or pustules should be assessed promptly
    \item Any sign of serious scalp infection with fever
\end{itemize}

\section*{Diagnosis and Investigations}
Diagnosis is usually clinical, with tests for atypical presentations.

\begin{itemize}
    \item \textbf{History and examination:} the pattern and onset of loss, family history, and scalp examination
    \item \textbf{Pull test:} assesses active shedding
    \item \textbf{Blood tests:} thyroid function, iron studies, and testosterone where indicated
    \item \textbf{Scalp biopsy:} for uncertain or scarring causes
    \item \textbf{Microscopy:} for fungal infection
    \item \textbf{Referral:} to a dermatologist for complex cases
\end{itemize}

\section*{Management}
Treatment depends on the cause and how much the hair loss bothers the person.

\begin{itemize}
    \item \textbf{Minoxidil lotion:} applied to the scalp, slows loss and can regrow some hair in pattern baldness; available over the counter
    \item \textbf{Finasteride tablets:} a prescription medicine that blocks DHT production, slowing loss and promoting regrowth in many men
    \item \textbf{Dutasteride:} a related medicine sometimes used
    \item \textbf{Low-level laser therapy:} a device-based option
    \item \textbf{Hair transplant:} surgery for permanent restoration
    \item \textbf{Alopecia areata:} steroid injections and other treatments
    \item \textbf{Telogen effluvium:} reassurance, as it resolves when the trigger passes
    \item \textbf{Treating underlying causes:} thyroid, iron, and scalp infections
    \item \textbf{Cosmetic options:} scalp micropigmentation, wigs, and concealers
\end{itemize}

\section*{Complications}
\begin{itemize}
    \item Psychological distress, reduced confidence, and effects on social and work life
    \item Permanent hair loss in untreated progressive pattern baldness
    \item Side effects of finasteride, including rare sexual side effects
    \item Scalp irritation from minoxidil
    \item Unrealistic expectations leading to dissatisfaction with treatment
    \item Missed diagnosis of underlying medical conditions
\end{itemize}

\section*{Prognosis and Outcomes}
Male pattern hair loss is progressive, but treatment can slow it substantially and regrow some hair, with the best results when started early. Minoxidil and finasteride are effective in many men and must be continued to maintain benefit. Alopecia areata often improves spontaneously or with treatment. Telogen effluvium resolves fully. For men who want a permanent solution, hair transplant surgery produces lasting results. The key is starting treatment early and having realistic expectations.

\section*{Prevention and Self-Care}
\begin{itemize}
    \item Start treatment early for pattern hair loss, as it preserves rather than regrows
    \item Use minoxidil and finasteride as directed and be consistent
    \item Be patient: results take four to six months to appear
    \item Eat a balanced diet with adequate iron and protein
    \item Avoid tight hairstyles and harsh chemical treatments
    \item Manage stress
    \item See a doctor for sudden, patchy, or unusual hair loss
    \item Seek support if hair loss is affecting your confidence
    \item Choose hair restoration professionals carefully if considering surgery
\end{itemize}

\section*{Sources and Further Reading}
The following reputable sources provide further information for healthcare students and patients seeking reliable, current advice about hair loss in men.

\begin{itemize}
    \item Healthdirect Australia. \textit{Male pattern hair loss.} https://www.healthdirect.gov.au/
    \item Australasian College of Dermatologists. \textit{Hair loss in men.} https://www.dermcoll.edu.au/
    \item Andrology Australia. \textit{Male hair loss information.} https://www.andrologyaustralia.org/
    \item NPS MedicineWise. \textit{Minoxidil and finasteride for hair loss.} https://www.nps.org.au/
    \item NHS. \textit{Male pattern baldness.} https://www.nhs.uk/
    \item American Academy of Dermatology. \textit{Hair loss in men.} https://www.aad.org/
\end{itemize}
